\documentclass[12pt, titlepage]{article}

\usepackage{amsmath, mathtools}

\usepackage[round]{natbib}
\usepackage{amsfonts}
\usepackage{amssymb}
\usepackage{graphicx}
\usepackage{colortbl}
\usepackage{xr}
\usepackage{hyperref}
\usepackage{longtable}
\usepackage{xfrac}
\usepackage{tabularx}
\usepackage{float}
\usepackage{siunitx}
\usepackage{booktabs}
\usepackage{multirow}
\usepackage[section]{placeins}
\usepackage{caption}
\usepackage{fullpage}

\hypersetup{
bookmarks=true,     % show bookmarks bar?
colorlinks=true,       % false: boxed links; true: colored links
linkcolor=red,          % color of internal links (change box color with linkbordercolor)
citecolor=blue,      % color of links to bibliography
filecolor=magenta,  % color of file links
urlcolor=cyan          % color of external links
}

\usepackage{array}

\externaldocument{../../SRS/SRS}

\input{../../Comments.text}
\input{../../Common.text}

\begin{document}

\title{Module Interface Specification for \progname{}}

\author{\authname}

\date{\today}

\maketitle

\pagenumbering{roman}

\section{Revision History}

\begin{tabularx}{\textwidth}{p{3cm}p{2cm}X}
\toprule {\bf Date} & {\bf Version} & {\bf Notes}\\
\midrule
Date 1 & 1.0 & Notes\\
Date 2 & 1.1 & Notes\\
\bottomrule
\end{tabularx}

~\newpage

\section{Symbols, Abbreviations and Acronyms}

See the SRS at
\url{https://github.com/capstoneCEGJM/Lake/blob/main/pdfs/SRS/SRS.pdf} for the
full list. The abbreviations and acronyms used in this document are:

\begin{center}
\renewcommand{\arraystretch}{1.2}
\begin{tabular}{l p{11cm}}
\toprule
\textbf{Symbol} & \textbf{Description} \\
\midrule
API & Application programming interface \\
GGML & The tensor library \progname{} uses to run model layers \\
GGUF & The file format that stores a model's weights, tokeniser and metadata \\
GPU & Graphics processing unit \\
I/O & Input/output \\
ID & Identifier \\
MG & Module Guide \\
MIS & Module Interface Specification \\
SRS & Software Requirements Specification \\
\bottomrule
\end{tabular}
\end{center}

\newpage

\tableofcontents

\newpage

\pagenumbering{arabic}

\section{Introduction}

The following document details the Module Interface Specifications for
\progname{}, a system that runs one large language model across several
devices on a local network. Each device holds a consecutive range of the
model's layers, and the devices pass hidden states to one another to generate
each token.

Complementary documents include the System Requirement Specifications
and Module Guide.  The full documentation and implementation can be
found at \url{https://github.com/capstoneCEGJM/Lake}.

\section{Notation}

The structure of the MIS for modules comes from \citet{HoffmanAndStrooper1995},
with the addition that template modules have been adapted from
\cite{GhezziEtAl2003}.  The mathematical notation comes from Chapter 3 of
\citet{HoffmanAndStrooper1995}.  For instance, the symbol := is used for a
multiple assignment statement and conditional rules follow the form $(c_1
\Rightarrow r_1 | c_2 \Rightarrow r_2 | ... | c_n \Rightarrow r_n )$.

The following table summarizes the primitive data types used by \progname. 

\begin{center}
\renewcommand{\arraystretch}{1.2}
\noindent 
\begin{tabular}{l l p{7.5cm}} 
\toprule 
\textbf{Data Type} & \textbf{Notation} & \textbf{Description}\\ 
\midrule
Boolean & $\mathbb{B}$ & \textit{true} or \textit{false}\\
character & char & a single symbol or digit\\
integer & $\mathbb{Z}$ & a number without a fractional component in (-$\infty$, $\infty$) \\
natural number & $\mathbb{N}$ & a number without a fractional component in [0, $\infty$) \\
real & $\mathbb{R}$ & any number in (-$\infty$, $\infty$)\\
\bottomrule
\end{tabular}
\end{center}

\noindent
The specification of \progname \ uses some derived data types:
\begin{itemize}
  \item \textbf{Strings} are sequences of characters, written string.
  \item \textbf{Sequences} hold elements of one type in order, written
    ``seq of $T$''. $\langle\rangle$ is the empty sequence and $|s|$ is the
    length of $s$.
  \item \textbf{Sets} hold elements of one type with no order or repeats,
    written ``set of $T$''. $\{\}$ is the empty set.
  \item \textbf{Tuples} hold named fields of possibly different types, written
    $(f_1: T_1, \dots, f_n: T_n)$. Field $f$ of tuple $r$ is $r.f$.
  \item \textbf{Maps} $A \to B$ are finite functions from keys of type $A$ to
    values of type $B$. $\mathrm{dom}(f)$ is the set of keys of $f$,
    $\{a \mapsto b\}$ is the map with one entry, and $f(a) := b$ changes or adds
    the entry for $a$.
  \item \textbf{Enumerated types} are written as the set of their values in
    sans-serif type, such as $\{\mathsf{idle}, \mathsf{member}\}$.
\end{itemize}

\noindent
Symbols taken from the SRS keep their meaning there. In particular,
partitioning schemes, the allocation states $\mathsf{U}$, $\mathsf{A}$,
$\mathsf{I}$, $\mathsf{G}$, $\mathsf{R}$ and $\mathsf{X}$, and the unit
memories $m_\ell$ come from SRS S.9. Symbolic constants in italics, such as
\textit{HEALTH\_INTERVAL}, have the values given in the SRS Symbolic Constants
appendix. Times are in seconds and memory sizes are in bytes.

\section{Module Decomposition}

The following table is taken directly from the Module Guide document for this project.

\begin{table}[h!]
\centering
\begin{tabular}{p{0.3\textwidth} p{0.6\textwidth}}
\toprule
\textbf{Level 1} & \textbf{Level 2}\\
\midrule

{Hardware-Hiding} & GGML\\
\midrule

\multirow{4}{0.3\textwidth}{Behaviour-Hiding} & Coordinator\\
& Core\\
& Inference Engine\\
& Lake Client\\
\midrule

\multirow{2}{0.3\textwidth}{Software Decision} & Partition\\
& Tensor I/O\\
\bottomrule

\end{tabular}
\caption{Module Hierarchy}
\label{TblMH}
\end{table}

\newpage
~\newpage

\section{MIS of Coordinator} \label{mis:coordinator}

\subsection{Module}

Coordinator

\subsection{Syntax}

\subsubsection{Exported Types}

\begin{itemize}
  \item $\mathit{DeviceID}$ and $\mathit{Report}$, as in the Core module
    (Section~\ref{mis:core}).
  \item $\mathit{DevState} = \{\mathsf{discovered}, \mathsf{invited},
    \mathsf{declined}, \mathsf{confirmed}, \mathsf{failed}\}$
  \item $\mathit{NetState} = \{\mathsf{U}, \mathsf{A}, \mathsf{I}, \mathsf{G},
    \mathsf{R}, \mathsf{X}\}$, the states of the allocation state machine in SRS
    S.9.3.
  \item $\mathit{Scheme} = \text{seq of } (\mathit{node}: \mathit{DeviceID},
    \mathit{first}: \mathbb{N}, \mathit{last}: \mathbb{N})$: a split of a model
    across nodes. Each entry gives a node and the first and last layers it
    holds, in the order the layers run.
  \item $\mathit{Request} = (\mathit{messages}: \text{seq of string},
    \mathit{maxTokens}: \mathbb{N}, \mathit{temperature}: \mathbb{R})$: a chat
    request from an application: the conversation so far,
    the most tokens to generate, and how random the sampling is, where 0 always
    picks the most likely token.
  \item $\mathit{Status} = (q: \mathit{NetState},$
    $\mathit{nodes}: \text{seq of } (\mathit{id}: \mathit{DeviceID},$
    $\mathit{state}: \mathit{DevState}, \mathit{report}: \mathit{Report}),$
    $P: \mathit{Scheme})$: a snapshot of a
    Lake network: its state, each known device with its progress and latest
    report, and the split in use.
\end{itemize}

\subsubsection{Exported Access Programs}

\begin{center}
\begin{tabular}{p{3.3cm} p{3.6cm} p{2.9cm} p{3.6cm}}
\toprule
\textbf{Name} & \textbf{In} & \textbf{Out} & \textbf{Exceptions} \\
\midrule
init & $\mathit{self}: \mathit{DeviceID}$ & -- & -- \\
inviteDevices & $S$: set of $\mathit{DeviceID}$ & -- & UnknownDevice \\
loadModel & $\mathit{file}$: string & -- & Unconfirmed, NoValidScheme \\
monitorNetwork & -- & -- & NotInitialised \\
serveRequest & $x: \mathit{Request}$ & seq of string & NotLoaded, RequestFailed \\
unloadModel & -- & -- & NotInitialised \\
getStatus & -- & $\mathit{Status}$ & NotInitialised \\
\bottomrule
\end{tabular}
\end{center}

{\raggedright \noindent inviteDevices, loadModel, unloadModel and getStatus
are called by the Lake Client. serveRequest is called once for each request
that an application sends to the OpenAI-compatible API.\par}

\subsection{Semantics}

\subsubsection{State Variables}

\begin{itemize}
  \item $\mathit{initialised}: \mathbb{B}$
  \item $\mathit{self}: \mathit{DeviceID}$, the leader node
  \item $\mathit{dev}: \mathit{DeviceID} \to \mathit{DevState}$
  \item $\mathit{rep}: \mathit{DeviceID} \to \mathit{Report}$, the latest
    report from each node
  \item $\mathit{lastSeen}: \mathit{DeviceID} \to \mathbb{R}$, the time on the
    leader's clock at which the latest report from each node arrived
  \item $\mathit{file}$: string, the GGUF file of the loaded model
  \item $P: \mathit{Scheme}$, the scheme the model is loaded on
  \item $P': \mathit{Scheme}$, the scheme being loaded
  \item $q: \mathit{NetState}$
  \item $\mathit{active}: \mathbb{N}$, the number of requests in progress
\end{itemize}

\subsubsection{Environment Variables}

\begin{itemize}
  \item $t: \mathbb{R}$, the current time in seconds.
  \item The local network, through which the Coordinator reaches the Core
    module on every node.
\end{itemize}

\subsubsection{Exceptions}

\begin{itemize}
  \item NotInitialised: any access program other than init is called before
    init. Every access program other than init can raise it.
  \item UnknownDevice: inviteDevices names a device that has not been discovered.
  \item Unconfirmed: loadModel is called while an invited device has not answered.
  \item NoValidScheme: no valid partitioning scheme exists for the Lake network.
  \item NotLoaded: serveRequest is called when no model is loaded.
  \item RequestFailed: a request still fails after \textit{RECOVERY\_LIMIT}
    retries.
\end{itemize}

\newpage

\section{MIS of Core} \label{mis:core}

\subsection{Module}

Core

\subsection{Syntax}

\subsubsection{Exported Types}

\begin{itemize}
  \item $\mathit{DeviceID}$: string. An identifier that is unique to each device
    on the local network.
  \item $\mathit{Report} = (\mathit{mem}: \mathbb{N}, \mathit{rate}:
    \mathbb{R})$: free GPU-accessible memory in bytes and token throughput in
    tokens per second.
  \item $\mathit{MemberState} = \{\mathsf{idle}, \mathsf{invited},
    \mathsf{member}, \mathsf{loaded}\}$: whether the node has no Lake network,
    has an invitation waiting for its owner's answer, has joined a Lake network,
    or holds a partition.
  \item $\mathit{CoreStatus} = (s: \mathit{MemberState}, \mathit{leader}:
    \mathit{DeviceID}, a: \mathbb{N}, b: \mathbb{N})$: the node's membership,
    its leader, and the first and last partition units it holds.
\end{itemize}

\subsubsection{Exported Access Programs}

\begin{center}
\begin{tabular}{p{3.3cm} p{4.3cm} p{2.0cm} p{3.8cm}}
\toprule
\textbf{Name} & \textbf{In} & \textbf{Out} & \textbf{Exceptions} \\
\midrule
init & $\mathit{self}: \mathit{DeviceID}$ & -- & -- \\
answerScan & $\mathit{from}: \mathit{DeviceID}$ & $\mathit{DeviceID}$ & NotInitialised \\
receiveInvite & $\mathit{from}: \mathit{DeviceID}$ & -- & Busy \\
answerInvite & $\mathit{ok}: \mathbb{B}$ & -- & NoInvitation \\
loadPartition & $\mathit{file}$: string, $a: \mathbb{N}$, $b: \mathbb{N}$,
  $\mathit{prev}: \mathit{DeviceID}$, $\mathit{next}: \mathit{DeviceID}$
  & -- & NotLeader, LoadFailed \\
releasePartition & -- & -- & NotLeader \\
reportHealth & -- & $\mathit{Report}$ & NotLeader \\
checkLeader & -- & -- & NotInitialised \\
getStatus & -- & $\mathit{CoreStatus}$ & NotInitialised \\
\bottomrule
\end{tabular}
\end{center}

\noindent answerScan, receiveInvite, loadPartition, releasePartition and
reportHealth are called by the leader's Coordinator. answerScan answers a
discovery scan from device $\mathit{from}$, and loadPartition returns once the partition is
loaded. answerInvite is called by the Lake Client with the device owner's
choice. checkLeader is called periodically to detect a leader that has
not queried the node for \textit{LEADER\_TIMEOUT}. The leader node
runs Core like every other node.

\subsection{Semantics}

\subsubsection{State Variables}

\begin{itemize}
  \item $\mathit{initialised}: \mathbb{B}$
  \item $\mathit{self}: \mathit{DeviceID}$, this node
  \item $s: \mathit{MemberState}$
  \item $\mathit{leader}: \mathit{DeviceID}$, the leader that invited this node
  \item $\mathit{file}$: string, the GGUF file the node reads its partition
    from
  \item $a, b: \mathbb{N}$, the first and last partition units the node holds
  \item $\mathit{prev}, \mathit{next}: \mathit{DeviceID}$, the nodes that send
    to and receive from this node
  \item $\mathit{lastQuery}: \mathbb{R}$, the time of the leader's last query
\end{itemize}

\subsubsection{Environment Variables}

\begin{itemize}
  \item $t: \mathbb{R}$, the current time in seconds.
  \item The local network, through which the Core module answers the leader's
    Coordinator.
  \item The device's GPU-accessible memory, whose free size is reported to the
    leader.
\end{itemize}

\subsubsection{Exceptions}

\begin{itemize}
  \item NotInitialised: any access program other than init is called before
    init. Every access program other than init can raise it.
  \item Busy: receiveInvite is called while the node is a member of a Lake network led
    by another device.
  \item NoInvitation: answerInvite is called when no invitation is waiting.
  \item NotLeader: loadPartition, releasePartition or reportHealth is called by from a device other than the
    node's leader, or before the node has joined.
  \item LoadFailed: the assigned partition cannot be loaded into the node's
    GPU-accessible memory.
\end{itemize}

\newpage

\section{MIS of \wss{Module Name}} \label{Module} \wss{Use labels for
  cross-referencing}

\wss{You can reference SRS labels, such as R\ref{R_Inputs}.}

\wss{It is also possible to use \LaTeX for hypperlinks to external documents.}

\subsection{Module}

\wss{Short name for the module}

\subsection{Uses}


\subsection{Syntax}

\subsubsection{Exported Constants}

\subsubsection{Exported Access Programs}

\begin{center}
\begin{tabular}{p{2cm} p{4cm} p{4cm} p{2cm}}
\hline
\textbf{Name} & \textbf{In} & \textbf{Out} & \textbf{Exceptions} \\
\hline
\wss{accessProg} & - & - & - \\
\hline
\end{tabular}
\end{center}

\subsection{Semantics}

\subsubsection{State Variables}

\wss{Not all modules will have state variables.  State variables give the module
  a memory.}

\subsubsection{Environment Variables}

\wss{This section is not necessary for all modules.  Its purpose is to capture
  when the module has external interaction with the environment, such as for a
  device driver, screen interface, keyboard, file, etc.}

\subsubsection{Assumptions}

\wss{Try to minimize assumptions and anticipate programmer errors via
  exceptions, but for practical purposes assumptions are sometimes appropriate.}

\subsubsection{Access Routine Semantics}

\noindent \wss{accessProg}():
\begin{itemize}
\item transition: \wss{if appropriate} 
\item output: \wss{if appropriate} 
\item exception: \wss{if appropriate} 
\end{itemize}

\wss{A module without environment variables or state variables is unlikely to
  have a state transition.  In this case a state transition can only occur if
  the module is changing the state of another module.}

\wss{Modules rarely have both a transition and an output.  In most cases you
  will have one or the other.}

\subsubsection{Local Functions}

\wss{As appropriate} \wss{These functions are for the purpose of specification.
  They are not necessarily something that is going to be implemented
  explicitly.  Even if they are implemented, they are not exported; they only
  have local scope.}

\newpage

\bibliographystyle {plainnat}
\bibliography {../../../refs/References}

\newpage

\newpage{}

\section{Appendix --- Generative AI Use Disclosure}

\wss{ChatGPT (version 5) was used to brainstorm ideas for the template for this disclose section.}

\subsection{AI Tools Used}

\begin{itemize}
    \item Anthropic Claude Opus 5.5, used through Claude Code.
\end{itemize}

\wss{Give the name of the AI tool(s) used and the version.}

\subsection{How and Where Used}

\begin{itemize}
    \item \textbf{MIS of Coordinator and MIS of Core}: Claude drafted the
    exported types, access programs, state variables, environment variables
    and exceptions from the SRS, the team's design decisions and the module
    split we chose. We reviewed the draft and directed changes to the names,
    types and access programs. Claude also drafted the Introduction, the
    Notation section, the Module Decomposition table and the Symbols,
    Abbreviations and Acronyms table.
\end{itemize}

\wss{Explain what the tool(s) were used for.  Examples include: brainstorming,
explaining a technical concept, reviewing the document, generating or modifying
text, generating or modifying diagrams, identifying errors or omissions or
inconsistencies.}

\wss{For each of the uses, identify which parts of the document were affected.}

\wss{You do not need to report every interaction, but you should disclose the 
uses that materially contributed to your work}

\subsection{Verification of AI-Generated Content}

The access programs of the Coordinator and Core were checked against each other
so that every call the Coordinator makes on Core exists in Core's
specification, and every type used is defined. The symbolic constants and
state names were checked against the SRS.

\wss{\begin{itemize}
    \item Describe how the team checked AI-generated or AI-assisted material
    before including it in the document.
    \item In particular, verify factual claims, technical assertions,
    requirements, calculations, references, and other information for which an
    AI system may produce incorrect or fabricated results.
    \item \textbf{The team is responsible for the accuracy, completeness,
consistency, and appropriateness of the submitted document, regardless of
whether material was generated by a human or an AI system.}
\end{itemize}}

\newpage{}

\section{Appendix --- Reflection}

\wss{Not required for CAS 741 projects}

The information in this section will be used to evaluate the team members on the
graduate attribute of Problem Analysis and Design.

\input{../../Reflection.text}

\begin{enumerate}
  \item What went well while writing this deliverable? 
  \item What pain points did you experience during this deliverable, and how
    did you resolve them?
  \item Which of your design decisions stemmed from speaking to your client(s)
  or a proxy (e.g. your peers, stakeholders, potential users)? For those that
  were not, why, and where did they come from?
  \item While creating the design doc, what parts of your other documents (e.g.
  requirements, hazard analysis, etc), it any, needed to be changed, and why?
  \item What are the limitations of your solution?  Put another way, given
  unlimited resources, what could you do to make the project better? (LO\_ProbSolutions)
  \item Give a brief overview of other design solutions you considered.  What
  are the benefits and tradeoffs of those other designs compared with the chosen
  design?  From all the potential options, why did you select the documented design?
  (LO\_Explores)
  \item (After you have implemented another team's module, which means this
  isn't filled in until after the original deadline). What did you learn by
  implementing another team's module? Were all the details you needed in the
  documentation, or did you need to make assumptions, or ask the other team
  questions?  If your team also had another team implement one of your modules,
  what was this experience like?  Are there things in your documentation you
  could have changed to make the process go more smoothly for when an
  ``outsider'' completes some of the implementation?

  Your reflection should be about concepts learned in your SE education, such as modular decomposition, cohesion, coupling, information hiding, etc.  You should critique the quality and utility of the specification you received.

  You should create a PR for your implemented module.  A link to the PR should be part of your reflection.

\end{enumerate}

\end{document}